%% thesis-chapter.tex
%% Template một chapter điển hình của luận văn thạc sĩ / tiến sĩ
%% Engine: XeLaTeX (cho tiếng Việt)
%%
%% File này thường được \include vào file main.tex chính:
%%   \include{chapters/chapter2-theory}
%%
%% Để compile standalone, bật \documentclass bên dưới.
%% Khi dùng cùng main.tex, xoá phần documentclass và begin/end document.

\documentclass[12pt, a4paper]{report}

%% ================================================================
%% PACKAGES (chỉ cần khi compile standalone)
%% ================================================================
\usepackage{fontspec}
\usepackage{polyglossia}
\setmainlanguage{vietnamese}
\setmainfont{Times New Roman}

\usepackage{geometry}
\geometry{a4paper, top=3cm, bottom=2.5cm, left=3.5cm, right=2cm}

\usepackage{graphicx}
\usepackage{float}
\usepackage{booktabs}
\usepackage{tabularx}
\usepackage{multirow}
\usepackage{amsmath}
\usepackage{amssymb}
\usepackage{hyperref}
\usepackage[backend=biber, style=numeric, sorting=none]{biblatex}
\addbibresource{sample.bib}
\usepackage[colorinlistoftodos]{todonotes}
\usepackage{setspace}
\onehalfspacing

%% Đánh số từ 1 cho standalone
\setcounter{chapter}{1}

\begin{document}

%% ================================================================
%% CHƯƠNG 2: CƠ SỞ LÝ THUYẾT
%% (Trong luận văn thực, đây là \include{chapters/chapter2})
%% ================================================================

\chapter{Cơ sở lý thuyết}
\label{chap:theory}

Chương này trình bày nền tảng lý thuyết liên quan đến bài toán phân loại
văn bản tiếng Việt, bao gồm các khái niệm về học máy, mạng nơ-ron sâu,
và các mô hình ngôn ngữ hiện đại.

%% ================================================================
\section{Học máy và Học sâu}
\label{sec:ml-dl}
%% ================================================================

\subsection{Học máy}

Học máy (\textit{Machine Learning --- ML}) là lĩnh vực nghiên cứu cách
máy tính tự học từ dữ liệu mà không cần lập trình tường minh từng quy tắc
\cite{Goodfellow2016}. Một bài toán học máy được định nghĩa bởi ba thành phần:

\begin{itemize}
  \item \textbf{Task} $T$: Bài toán cần giải quyết (phân loại, hồi quy, sinh ảnh...)
  \item \textbf{Experience} $E$: Dữ liệu huấn luyện
  \item \textbf{Performance measure} $P$: Chỉ số đánh giá hiệu suất
\end{itemize}

\subsection{Học sâu}

Học sâu (\textit{Deep Learning --- DL}) là một nhánh của học máy sử dụng
mạng nơ-ron nhân tạo nhiều lớp (\textit{deep neural networks}) để học
các biểu diễn phân cấp từ dữ liệu thô \cite{LeCun2015deep}.

\begin{equation}
  \label{eq:forward}
  h^{(l)} = \sigma\!\left(W^{(l)} h^{(l-1)} + b^{(l)}\right)
\end{equation}

trong đó $h^{(l)}$ là activation của lớp $l$, $W^{(l)}$ và $b^{(l)}$
là trọng số và bias tương ứng, $\sigma(\cdot)$ là hàm kích hoạt phi tuyến.

Hàm mất mát cross-entropy thường dùng cho bài toán phân loại nhiều lớp:

\begin{equation}
  \label{eq:crossentropy}
  \mathcal{L} = -\frac{1}{N}\sum_{i=1}^{N}\sum_{c=1}^{C}
  y_{ic} \log \hat{y}_{ic}
\end{equation}

trong đó $N$ là số mẫu, $C$ là số lớp, $y_{ic} \in \{0,1\}$ là nhãn thực,
và $\hat{y}_{ic} \in [0,1]$ là xác suất dự đoán.

%% ================================================================
\section{Kiến trúc Transformer}
\label{sec:transformer}
%% ================================================================

\subsection{Tổng quan}

Transformer \cite{Vaswani2017attention} là kiến trúc mạng nơ-ron
dựa hoàn toàn trên cơ chế attention, loại bỏ hoàn toàn thành phần
hồi quy (RNN) và tích chập (CNN). Điều này cho phép song song hóa
hoàn toàn trong quá trình huấn luyện.

\subsection{Scaled Dot-Product Attention}

Cơ chế attention cơ bản được tính như Phương trình~\eqref{eq:attention}:

\begin{equation}
  \label{eq:attention}
  \text{Attention}(Q, K, V) = \text{softmax}\!\left(\frac{QK^\top}{\sqrt{d_k}}\right)V
\end{equation}

trong đó:
\begin{itemize}
  \item $Q \in \mathbb{R}^{n \times d_k}$: Ma trận query
  \item $K \in \mathbb{R}^{m \times d_k}$: Ma trận key
  \item $V \in \mathbb{R}^{m \times d_v}$: Ma trận value
  \item $\sqrt{d_k}$: Hệ số scale để ổn định gradient
\end{itemize}

\subsection{Multi-Head Attention}

Multi-head attention cho phép mô hình chú ý đến thông tin từ nhiều
không gian biểu diễn khác nhau đồng thời:

\begin{align}
  \text{MultiHead}(Q,K,V) &= \text{Concat}(\text{head}_1, \ldots, \text{head}_h)\,W^O \\
  \text{where head}_i     &= \text{Attention}(QW_i^Q, KW_i^K, VW_i^V)
\end{align}

Hình~\ref{fig:transformer} minh họa kiến trúc Transformer encoder.

\begin{figure}[htbp]
  \centering
  %% Thay bằng file ảnh thực tế khi có
  \fbox{\parbox{0.7\textwidth}{\centering\vspace{1cm}
    [Hình minh họa kiến trúc Transformer]\\[0.3cm]
    \small Thay thế bằng \texttt{figures/transformer.png}
    \vspace{1cm}}}
  \caption{Kiến trúc Transformer encoder với multi-head self-attention
           và feed-forward network~\cite{Vaswani2017attention}.}
  \label{fig:transformer}
\end{figure}

%% ================================================================
\section{Các mô hình ngôn ngữ tiền huấn luyện}
\label{sec:plm}
%% ================================================================

\subsection{BERT}

BERT (\textit{Bidirectional Encoder Representations from Transformers})
\cite{Vaswani2017attention} là mô hình ngôn ngữ được tiền huấn luyện
trên hai nhiệm vụ tự giám sát:

\begin{enumerate}
  \item \textbf{Masked Language Modeling (MLM):} Dự đoán 15\% token
        bị che ngẫu nhiên trong câu
  \item \textbf{Next Sentence Prediction (NSP):} Dự đoán xem hai câu
        có liên tiếp nhau không
\end{enumerate}

\subsection{PhoBERT}

PhoBERT \cite{Le2024bert} là mô hình BERT được tiền huấn luyện đặc biệt
cho tiếng Việt, với các đặc điểm nổi bật:

\begin{table}[htbp]
  \centering
  \caption{Thông số kỹ thuật của PhoBERT-base và PhoBERT-large}
  \label{tab:phobert}
  \begin{tabular}{lcc}
    \toprule
    Tham số & PhoBERT-base & PhoBERT-large \\
    \midrule
    Số lớp Transformer        & 12     & 24 \\
    Kích thước hidden          & 768    & 1024 \\
    Số attention heads         & 12     & 16 \\
    Số tham số                 & 135M   & 370M \\
    Dữ liệu tiền huấn luyện    & 20GB   & 20GB \\
    Từ vựng                    & 64.000 & 64.000 \\
    \bottomrule
  \end{tabular}
\end{table}

Bảng~\ref{tab:phobert} cho thấy PhoBERT-large có dung lượng lớn hơn
gần 3 lần so với PhoBERT-base, tuy nhiên điều này không đảm bảo kết quả
tốt hơn trên mọi bài toán --- đặc biệt khi dữ liệu fine-tuning nhỏ.
\todo{Bổ sung kết quả so sánh PhoBERT-base vs large trên VNews-50K}

%% ================================================================
\section{Đánh giá mô hình}
\label{sec:evaluation}
%% ================================================================

\subsection{Các chỉ số đánh giá}

Trong bài toán phân loại đa lớp, chúng tôi sử dụng các chỉ số sau:

\begin{description}
  \item[Accuracy] Tỷ lệ mẫu được phân loại đúng trên tổng số mẫu:
        $\text{Acc} = \frac{\text{TP} + \text{TN}}{N}$

  \item[Precision] Trong các mẫu được dự đoán là lớp $c$, bao nhiêu
        phần trăm thực sự thuộc lớp $c$:
        $P = \frac{\text{TP}}{\text{TP} + \text{FP}}$

  \item[Recall] Trong các mẫu thực sự thuộc lớp $c$, bao nhiêu phần
        trăm được nhận diện đúng:
        $R = \frac{\text{TP}}{\text{TP} + \text{FN}}$

  \item[F1-score] Trung bình điều hòa của Precision và Recall:
        \begin{equation}
          \label{eq:f1}
          F_1 = 2 \cdot \frac{P \cdot R}{P + R}
        \end{equation}
\end{description}

Với bài toán đa lớp, chúng tôi sử dụng \textbf{macro-averaged F1},
tính trung bình F1 của từng lớp mà không trọng số theo số lượng mẫu:

\begin{equation}
  \label{eq:macro-f1}
  F_1^{\text{macro}} = \frac{1}{C} \sum_{c=1}^{C} F_{1,c}
\end{equation}

\subsection{Kiểm định thống kê}

Để đảm bảo kết quả so sánh có ý nghĩa thống kê, chúng tôi sử dụng
phép kiểm định McNemar \cite{Goodfellow2016} với ngưỡng $p < 0.05$.
Các thực nghiệm được lặp lại 5 lần với các random seed khác nhau,
kết quả báo cáo là mean $\pm$ std.

%% ================================================================
\section{Tóm tắt chương}
\label{sec:summary}
%% ================================================================

Chương này đã trình bày các nền tảng lý thuyết quan trọng:

\begin{enumerate}
  \item Học máy và học sâu với các kiến trúc mạng nơ-ron cơ bản
        (Phần~\ref{sec:ml-dl})
  \item Kiến trúc Transformer và cơ chế self-attention
        (Phần~\ref{sec:transformer})
  \item Các mô hình ngôn ngữ tiền huấn luyện, đặc biệt là PhoBERT
        phù hợp với tiếng Việt (Phần~\ref{sec:plm})
  \item Các chỉ số đánh giá cho bài toán phân loại đa lớp
        (Phần~\ref{sec:evaluation})
\end{enumerate}

Chương~\ref{chap:method} tiếp theo sẽ trình bày phương pháp đề xuất
dựa trên các nền tảng lý thuyết này.

%% ================================================================
%% BIBLIOGRAPHY (chỉ cần khi compile standalone)
%% ================================================================
\printbibliography[title={Tài liệu tham khảo}]

\end{document}
